% Claim proof file for row claim_3_20 -- "AcyclicNonCollidersBlockable".
% Auto-generated stub by scaffold/solve_chapter.py at row start. A manager
% agent dedicated to writing the TeX proof fills in the proof body below.
%
% This file is a sibling of `main.tex` inside `<subsection>/tex/`. The
% `subfiles` package lets it render both standalone and as part of
% `main.tex` via `\subfile{...}`.
%
% The statement is restated above the proof so the file is self-contained
% when rendered on its own. The proof must:
%   - Stay close to the lecture notes' proof if one exists (always search
%     the LN first for a `\begin{proof}` block immediately following the
%     claim; copy / lightly edit when present).
%   - Otherwise use the LN paradigm and cite prior claims by their `ref`.
%   - Be self-contained enough that a mathematician can follow it without
%     re-reading the LN.
%   - Have no `sorry`, `omitted`, or "obvious" shortcuts.

\documentclass[main]{subfiles}

\begin{document}
% ref: claim_3_20 (proof, with statement restated for readability)
% title: AcyclicNonCollidersBlockable
\def\rowref{claim\_3\_20}
\phantomsection\label{claim_3_20}

% --- statement (restated, synced verbatim from
% --- claim_3_20_statement_AcyclicNonCollidersBlockable.tex) ----------------
\begin{Rem}[AcyclicNonCollidersBlockable]
    Let $G = (J, V, E, L)$ be a CDMG (def \ref{def-cdmg}) and suppose $G$ is acyclic in the sense of def \ref{def-acylic}. Then, for every integer $n \ge 0$, every walk
    \[
        \pi = (v_0, a_0, v_1, a_1, \dots, v_{n-1}, a_{n-1}, v_n)
    \]
    in $G$ in the sense of def \ref{def:walks}, item~i.\ (so $v_0, \dots, v_n \in J \cup V$, $a_0, \dots, a_{n-1} \in E \cup L$, and the walk constraint
    \[
        \bigl(a_k = (v_k, v_{k+1}) \;\land\; a_k \in E \cup L\bigr)
        \;\lor\;
        \bigl(a_k = (v_{k+1}, v_k) \;\land\; a_k \in E\bigr)
    \]
    holds at every $k \in \{0, 1, \dots, n - 1\}$), and every position $k \in \{0, 1, \dots, n\}$, the following implication holds:
    \[
        \bigl(\text{$k$ is a non-collider on $\pi$ per def \ref{def:collider_noncollider}, item~i.}\bigr)
        \;\implies\;
        \bigl(\text{$k$ is a blockable non-collider on $\pi$ per def \ref{def:unblockable_noncollider}}\bigr).
    \]
\end{Rem}

% --- proof ------------------------------------------------------------------
\begin{proof}
% Guided by the LN author's commented-out sketch at
% lecture-notes/lecture_notes/graphs.tex:1599-1614, adapted to the
% def_3_16 framing of this formalization: the LN sketch's conditions
% "v_k \notin \Sc^G(v_{k \pm 1})" are rewritten as
% "v_{k \pm 1} \in \Sc^G(v_k)" per def_3_16, item iii. The walk-edge
% writing, the end-position branch, and the "concatenation yields a
% directed walk of length \ge 1" step are spelled out explicitly so the
% proof file is self-contained.
Fix $n \ge 0$, a walk $\pi = (v_0, a_0, v_1, a_1, \dots, v_{n-1}, a_{n-1}, v_n)$ in $G$ (def \ref{def:walks}, item~i.), and a position $k \in \{0, 1, \dots, n\}$. Suppose $k$ is a non-collider on $\pi$ (def \ref{def:collider_noncollider}, item~i.). We must show $k$ is a blockable non-collider on $\pi$ in the sense of def \ref{def:unblockable_noncollider}, namely:
\begin{itemize}
    \item[(B1)] $k$ is a non-collider on $\pi$ in the sense of def \ref{def:collider_noncollider}, item~i.; AND
    \item[(B2)] $k$ is not an unblockable non-collider on $\pi$ in the sense of def \ref{def:unblockable_noncollider}.
\end{itemize}
Clause (B1) is the hypothesis. It remains to establish clause (B2).

We split on the position $k$ along $\pi$.

\medskip\noindent\textbf{Case 1 -- end-position ($k = 0$ or $k = n$).}
By def \ref{def:unblockable_noncollider}, clause~ii., any unblockable non-collider on $\pi$ must be \emph{interior}, i.e.\ $k \notin \{0, n\}$ (equivalently $1 \le k \le n - 1$). The hypothesis $k \in \{0, n\}$ of this case directly negates clause~ii.\ of def \ref{def:unblockable_noncollider}. Therefore $k$ is not an unblockable non-collider on $\pi$, establishing (B2). Hence $k$ is a blockable non-collider on $\pi$ in this case.

\medskip\noindent\textbf{Case 2 -- interior position ($1 \le k \le n - 1$).}
We argue by contradiction: suppose, toward contradiction, that $k$ \emph{is} an unblockable non-collider on $\pi$ in the sense of def \ref{def:unblockable_noncollider}. We will derive a non-trivial directed walk from $v_k$ to itself in $G$, contradicting def \ref{def-acylic}.

\smallskip\noindent\emph{Step 2.1 -- the non-collider hypothesis yields an outgoing walk-edge of $v_k$.}
Since $k$ is a non-collider on $\pi$, def \ref{def:collider_noncollider}, item~i.\ gives $\mathrm{ah}_\pi(k) \le 1$, where
\[
    \mathrm{ah}_\pi(k) = \Bigl|\Bigl\{\, i \in \{k - 1, k\} \,\Bigm|\, 0 \le i \le n - 1 \,\land\, a_i \text{ is an edge into } v_k \text{ in } G \,\Bigr\}\Bigr|.
\]
Since $1 \le k \le n - 1$, both indices $i \in \{k - 1, k\}$ are admissible (they satisfy $0 \le i \le n - 1$), so the index set defining $\mathrm{ah}_\pi(k)$ is $\{k - 1, k\}$ and $|\{k - 1, k\}| = 2$. Hence $\mathrm{ah}_\pi(k) \le 1$ means at least one of the two walk-incident edges $a_{k - 1}, a_k$ is \emph{not} an edge into $v_k$ in the sense of def \ref{def-edge-relations}, item~ii.

We expand what ``not an edge into $v_k$'' forces on each walk-incident edge using the walk constraint of def \ref{def:walks}, item~i.\ (recalled in the statement above). For each admissible index $i \in \{k - 1, k\}$, the walk constraint at $i$ states
\[
    \bigl(a_i = (v_i, v_{i+1}) \;\land\; a_i \in E \cup L\bigr)
        \;\lor\;
    \bigl(a_i = (v_{i+1}, v_i) \;\land\; a_i \in E\bigr),
\]
so $a_i$ is one of three configurations:
\begin{itemize}
    \item[(C1)] $a_i = (v_i, v_{i+1}) \in E$ (a directed edge from $v_i$ to $v_{i+1}$);
    \item[(C2)] $a_i = (v_i, v_{i+1}) \in L$ (a bidirected edge with the ordered-pair writing $(v_i, v_{i+1})$);
    \item[(C3)] $a_i = (v_{i+1}, v_i) \in E$ (a directed edge from $v_{i+1}$ to $v_i$).
\end{itemize}
Applying the ``edge into $v_k$'' predicate of def \ref{def-edge-relations}, item~ii.\ ($a_i \in E$ with second coordinate $v_k$, or $a_i \in L$ with either coordinate $v_k$) to each configuration at slot $i \in \{k - 1, k\}$ (so the writing $a_i = (e_1, e_2)$ takes $e_1, e_2 \in \{v_i, v_{i+1}\}$ in some order, and $v_k \in \{v_i, v_{i+1}\}$ because $i \in \{k - 1, k\}$ means $\{v_i, v_{i+1}\}$ contains $v_k$):
\begin{itemize}
    \item For \emph{slot $i = k - 1$} ($v_i = v_{k - 1}, v_{i+1} = v_k$):
    \begin{itemize}
        \item (C1) $a_{k-1} = (v_{k-1}, v_k) \in E$: into $v_k$ (the $E$-clause fires, $e_2 = v_k$).
        \item (C2) $a_{k-1} = (v_{k-1}, v_k) \in L$: into $v_k$ (the $L$-clause fires, $e_2 = v_k$).
        \item (C3) $a_{k-1} = (v_k, v_{k-1}) \in E$: \emph{not} into $v_k$ (the $E$-clause requires $e_2 = v_k$, but $e_2 = v_{k-1}$; and $L$-clause is excluded since $a_{k-1} \in E$). This is the only configuration at slot $k - 1$ for which $a_{k - 1}$ fails the ``edge into $v_k$'' predicate.
    \end{itemize}
    \item For \emph{slot $i = k$} ($v_i = v_k, v_{i+1} = v_{k+1}$):
    \begin{itemize}
        \item (C1) $a_k = (v_k, v_{k+1}) \in E$: \emph{not} into $v_k$ (the $E$-clause requires $e_2 = v_k$, but $e_2 = v_{k+1}$; and $L$-clause is excluded since $a_k \in E$). This is the only configuration at slot $k$ for which $a_k$ fails the ``edge into $v_k$'' predicate.
        \item (C2) $a_k = (v_k, v_{k+1}) \in L$: into $v_k$ (the $L$-clause fires, $e_1 = v_k$).
        \item (C3) $a_k = (v_{k+1}, v_k) \in E$: into $v_k$ (the $E$-clause fires, $e_2 = v_k$).
    \end{itemize}
\end{itemize}
In summary, an admissible slot $i \in \{k - 1, k\}$ has $a_i$ \emph{not} an edge into $v_k$ \emph{iff} either
\[
    i = k - 1 \;\text{and}\; a_{k - 1} = (v_k, v_{k - 1}) \in E,
    \qquad\text{or}\qquad
    i = k \;\text{and}\; a_k = (v_k, v_{k + 1}) \in E.
\]
By the previous paragraph, the non-collider hypothesis forces at least one such slot to occur. We capture this in a single statement: there exists a slot $i^* \in \{k - 1, k\}$ at which $a_{i^*}$ is an \emph{outgoing walk-edge of $v_k$ at position $k$ on $\pi$} in the sense of def \ref{def:unblockable_noncollider} (the predicate ``$a_i \in E \land e_1 = v_k$''), with the two sub-cases:
\begin{itemize}
    \item[(O1)] $i^* = k - 1$ with $a_{k - 1} = (v_k, v_{k - 1}) \in E$ -- the other endpoint of $a_{k - 1}$ along $\pi$ is $v_{k - 1}$;
    \item[(O2)] $i^* = k$ with $a_k = (v_k, v_{k + 1}) \in E$ -- the other endpoint of $a_k$ along $\pi$ is $v_{k + 1}$.
\end{itemize}
(Both sub-cases may simultaneously hold -- corresponding to the ``fork'' pattern -- but for the argument below we just need at least one.)

\smallskip\noindent\emph{Step 2.2 -- the unblockable assumption forces the other endpoint into $\Sc^G(v_k)$.}
By the contradiction hypothesis, $k$ is unblockable on $\pi$, so clause~iii.\ of def \ref{def:unblockable_noncollider} fires:
\[
    \bigl(a_{k-1} = (v_k, v_{k-1}) \in E \;\implies\; v_{k-1} \in \Sc^G(v_k)\bigr)
        \;\land\;
    \bigl(a_k = (v_k, v_{k+1}) \in E \;\implies\; v_{k+1} \in \Sc^G(v_k)\bigr).
\]
Pick the slot $i^*$ from Step~2.1 and write $w$ for the other walk-endpoint at that slot:
\[
    w := v_{k - 1} \text{ if } i^* = k - 1, \qquad w := v_{k + 1} \text{ if } i^* = k.
\]
In either sub-case the antecedent of the corresponding conjunct fires (by the very definition of $i^*$ in Step~2.1), and we conclude
\[
    w \in \Sc^G(v_k),
\]
together with the directed-edge fact $(v_k, w) \in E$ recorded in Step~2.1.

\smallskip\noindent\emph{Step 2.3 -- $w \in \Sc^G(v_k)$ supplies a directed walk from $w$ to $v_k$.}
By def \ref{def_3_5}, item~vii.,
\[
    \Sc^G(v_k) \;=\; \Anc^G(v_k) \,\cap\, \Desc^G(v_k),
\]
so $w \in \Sc^G(v_k)$ entails $w \in \Anc^G(v_k)$. By def \ref{def_3_5}, item~iv., $w \in \Anc^G(v_k)$ means there exists a directed walk $\rho$ from $w$ to $v_k$ in $G$ in the sense of def \ref{def:walks}, item~ii.\ -- a tuple
\[
    \rho \;=\; (w = u_0, b_0, u_1, b_1, \dots, u_{m-1}, b_{m-1}, u_m = v_k)
\]
for some $m \ge 0$, with each $b_j = (u_j, u_{j+1}) \in E$ pointing forward along $\rho$ from $u_j$ to $u_{j+1}$.

\smallskip\noindent\emph{Step 2.4 -- concatenating the single edge $(v_k, w) \in E$ with $\rho$ yields a non-trivial directed walk from $v_k$ to itself.}
We assemble the directed walk
\[
    \widetilde{\rho} \;:=\; (v_k, \;(v_k, w), \;w = u_0, \;b_0, \;u_1, \;b_1, \;\dots, \;u_{m-1}, \;b_{m-1}, \;u_m = v_k).
\]
That is, prepend the node $v_k$ and the edge $(v_k, w) \in E$ to $\rho$. We verify that $\widetilde{\rho}$ is a directed walk from $v_k$ to itself in the sense of def \ref{def:walks}, item~ii.
\begin{itemize}
    \item \emph{Node and edge typing.} The nodes of $\widetilde{\rho}$ are $v_k, w = u_0, u_1, \dots, u_m = v_k$, all in $J \cup V$ ($v_k$ by hypothesis of the walk $\pi$, the $u_j$ by the directed-walk $\rho$). The edges are $(v_k, w)$ followed by $b_0, b_1, \dots, b_{m-1}$, all in $E$ ($(v_k, w) \in E$ from Step~2.1, the $b_j \in E$ from $\rho$).
    \item \emph{Length.} The length of $\widetilde{\rho}$, counted by the number of edges, is $1 + m$, where $m \ge 0$. Hence $1 + m \ge 1$, so $\widetilde{\rho}$ is a non-trivial walk (its length is at least~$1$, i.e.\ it has at least one edge).
    \item \emph{Endpoints.} The first node of $\widetilde{\rho}$ is $v_k$ (by construction), and the last node is $u_m = v_k$ (the terminus of $\rho$). So $\widetilde{\rho}$ starts and ends at $v_k$.
    \item \emph{Directed-walk shape.} Per def \ref{def:walks}, item~ii., a directed walk has all its edges of the form ``forward'' $(u_j, u_{j+1}) \in E$ at each successive slot, with arrowheads consistently pointing forward. The first edge of $\widetilde{\rho}$ is $(v_k, w) \in E$, fitting the slot writing $(\text{first node}, \text{second node}) = (v_k, w)$ in $E$; the subsequent edges $b_0, b_1, \dots, b_{m-1}$ are the directed-walk edges of $\rho$, hence fit the same forward-slot writing relative to $\rho$'s nodes; and the join between $(v_k, w)$ and $\rho$ at the node $w = u_0$ is consistent (the head of $(v_k, w)$ equals the start of $\rho$). All arrowheads point forward (from $v_k$ to $w$, then from $u_0$ through to $u_m$); no arrowhead points backward. Hence $\widetilde{\rho}$ has the directed-walk shape of def \ref{def:walks}, item~ii.
\end{itemize}
Therefore $\widetilde{\rho}$ is a directed walk from $v_k$ to $v_k$ in $G$ of length at least~$1$ -- equivalently, a \emph{non-trivial} directed walk from $v_k$ to itself in $G$.

\smallskip\noindent\emph{Step 2.5 -- contradiction with acyclicity.}
By def \ref{def-acylic}, $G$ acyclic asserts that no non-trivial directed walk from any node to itself exists in $G$. But Step~2.4 has constructed exactly such a walk $\widetilde{\rho}$ at the node $v_k \in J \cup V$. Contradiction.

\smallskip
Hence the assumption that $k$ is an unblockable non-collider on $\pi$ is false. So $k$ is not an unblockable non-collider on $\pi$, establishing (B2) in the interior case.

\medskip
In both cases (end-position and interior), the hypothesis ``$k$ is a non-collider on $\pi$'' implies $k$ is not an unblockable non-collider on $\pi$. Combining with (B1), $k$ is a blockable non-collider on $\pi$ in the sense of def \ref{def:unblockable_noncollider}.
\end{proof}

\end{document}
